\documentclass[11pt]{article}
\usepackage[a4paper,margin=1in]{geometry}
\usepackage{amsmath,amssymb,amsthm}
\usepackage[T1]{fontenc}
\usepackage{lmodern}
\usepackage[hidelinks]{hyperref}
\setlength{\emergencystretch}{3em}

\theoremstyle{plain}
\newtheorem{theorem}{Theorem}
\newtheorem{lemma}[theorem]{Lemma}
\theoremstyle{definition}
\newtheorem{definition}[theorem]{Definition}
\theoremstyle{remark}
\newtheorem{remark}[theorem]{Remark}

\newcommand{\Z}{\mathbb Z}
\newcommand{\Q}{\mathbb Q}
\newcommand{\End}{\operatorname{End}}

\title{A Counterexample to Conjecture 00000004280:\\
Isomorphic Endomorphism Rings Already in Rank $1$}
\author{%
  Feiyu\thanks{Independent researcher.}
  \and
  Claude (Opus 5.5)\thanks{Anthropic.}%
}
\date{2026-10-03}

\begin{document}
\maketitle

\begin{abstract}
Conjecture 00000004280 claims that there are non-isomorphic abelian groups with
isomorphic endomorphism rings, and that the minimal rank of such a pair is $3$. Let
$B\le\Q$ be the subgroup generated by the fractions $1/p$ with $p$ prime. Then $\Z$ and
$B$ both have rank $1$, their endomorphism rings are both isomorphic to $\Z$, and
$B\not\cong\Z$. So the minimal rank is at most $1$. The refutation is checked in
Lean~4 against Mathlib.
\end{abstract}

\section{The conjecture as stated}

The original text of conjecture 00000004280 reads:

\begin{quote}
Definition: The endomorphism ring is the ring of additive homomorphisms of a group to
itself. Conjecture: There exist two abelian groups whose endomorphism rings are
isomorphic as rings while the groups are not isomorphic, and the minimal rank of
realization is $3$.
\end{quote}

The Chinese version in \texttt{conjectures/00000004280.md} states the same two
claims.

\section{Reading of the statement}

A \emph{realization} is a pair $(A,B)$ of abelian groups with $\End(A)\cong\End(B)$ as
rings and $A\not\cong B$. The \emph{rank} of an abelian group $A$ is its torsion-free
rank, the rank of $A$ as a $\Z$-module (the dimension of $\Q\otimes A$). Our groups are
torsion-free, so this agrees with the total rank.

The statement does not say how the rank of a pair is measured: by the rank of the
first group, by the minimum, or by the maximum of the two ranks. Each reading implies
the following weakest form of the second claim:
\begin{equation}\label{eq:M}
\text{every realization $(A,B)$ has $\operatorname{rank}A\ge3$ or
$\operatorname{rank}B\ge3$.}\tag{M}
\end{equation}
We refute (M); the first claim (existence) is in fact confirmed by our example.

\section{Result}

Let $B=\langle 1/p : p\text{ prime}\rangle\le\Q$.

\begin{theorem}\label{thm:main}
$\End(\Z)\cong\Z\cong\End(B)$ as rings, $B\not\cong\Z$, and
$\operatorname{rank}\Z=1$, $\operatorname{rank}B\le1$. Hence $(\Z,B)$ is a
realization of rank at most $1$, (M) fails, and Conjecture 00000004280 is false.
\end{theorem}

\section{Proof}

\begin{lemma}\label{lem:val}
For every prime $p$ and every nonzero $x\in B$, the $p$-adic valuation satisfies
$v_p(x)\ge-1$.
\end{lemma}

\begin{proof}
The generators $1/q$ have $v_p(1/q)\in\{0,-1\}$. The set of $x$ with $x=0$ or
$v_p(x)\ge-1$ contains $0$ and is closed under negation and addition, because
$v_p(x+y)\ge\min(v_p(x),v_p(y))$.
\end{proof}

\begin{lemma}\label{lem:endo}
Every additive endomorphism $f$ of $B$ is multiplication by an integer.
\end{lemma}

\begin{proof}
First, $1=2\cdot\frac12\in B$. Let $x\in B$, written in lowest terms as $x=a/b$. Then
$b\cdot x=a\cdot1$ in $B$, so $b\,f(x)=a\,f(1)$ in $\Q$, i.e.\ $f(x)=x\,c$ with
$c=f(1)\in B$. For every prime $p$ we get $c/p=f(1/p)\in B$, so if $c\ne0$ then
$v_p(c)-1=v_p(c/p)\ge-1$ by Lemma~\ref{lem:val}, i.e.\ $v_p(c)\ge0$ for all $p$.
Hence $c$ is an integer: otherwise a prime $q$ divides its denominator, and $v_q(c)<0$.
\end{proof}

\begin{proof}[Proof of Theorem~\ref{thm:main}]
\emph{Endomorphism rings.} For any abelian group $G$, the canonical ring map
$\Z\to\End(G)$ sends $m$ to multiplication by $m$. For $G=\Z$ it is bijective, since an
endomorphism of $\Z$ is determined by $f(1)$. For $G=B$ it is surjective by
Lemma~\ref{lem:endo}, and injective because $m\cdot1=0$ in $B$ forces $m=0$. Hence
$\End(\Z)\cong\Z\cong\End(B)$.

\emph{Non-isomorphism.} Suppose $e\colon\Z\to B$ is an isomorphism and let
$k=e^{-1}(1)\ne0$. For every prime $p$ we have $p\cdot\frac1p=1$ in $B$, so
$k=p\,e^{-1}(1/p)$ is divisible by $p$. Taking a prime $p>|k|$ gives a contradiction.

\emph{Ranks.} $\operatorname{rank}\Z=1$. Any two elements $x=a/b$ and $y=c/d$ of $B$
(in lowest terms) satisfy $(cb)\,x-(ad)\,y=0$. If both are nonzero the coefficients
are nonzero, and if one is zero the pair is dependent anyway. So no two elements of $B$
are linearly independent over $\Z$, and $\operatorname{rank}B\le1$.
\end{proof}

\section{Formalization}

\sloppy
The Lean~4 project in \texttt{lean/} (file \texttt{lean/Submission/Basic.lean})
formalizes the definitions above and proves Theorem~\ref{thm:main} against Mathlib.
\begin{itemize}
  \item \texttt{Separates A B} is $\End(A)\cong\End(B)$ as rings (a
    \texttt{RingEquiv} between Mathlib's \texttt{AddMonoid.End A} and
    \texttt{AddMonoid.End B}) together with \texttt{IsEmpty (A} $\simeq_+$
    \texttt{B)}. \texttt{MinimalRankThree} is (M), with rank Mathlib's
    \texttt{Module.rank} over $\Z$.
  \item \texttt{Bsub} is $B$, defined as \texttt{AddSubgroup.closure} of
    $\{1/p\}$.
  \item \texttt{val\_ge} is Lemma~\ref{lem:val}, using Mathlib's
    \texttt{padicValRat}. \texttt{den\_eq\_one} and \texttt{endo\_apply} give
    Lemma~\ref{lem:endo}.
  \item \texttt{castZ\_bijective} and \texttt{castB\_bijective} show that
    \texttt{Int.castRingHom} into the two endomorphism rings is bijective, and
    \texttt{endEquiv} is the resulting ring isomorphism.
  \item \texttt{not\_iso} is the non-isomorphism, and \texttt{rank\_Z},
    \texttt{rank\_B\_le} are the rank computations.
  \item \texttt{conjecture\_00000004280\_false} states
    $\lnot(P\wedge\texttt{MinimalRankThree})$ for an arbitrary proposition $P$
    standing for the first claim.
\end{itemize}
The toolchain is \texttt{leanprover/lean4:v4.33.1}, Mathlib is pinned to
\texttt{v4.33.1}, and \texttt{lake build} succeeds. The project contains no
\texttt{sorry}, no \texttt{admit} and no \texttt{native\_decide}, and introduces no
axioms: \texttt{\#print axioms conjecture\_00000004280\_false} reports exactly
\texttt{propext}, \texttt{Classical.choice} and \texttt{Quot.sound}.

\fussy
\section{Remarks}

\begin{remark}
In the language of types of rank-$1$ groups, $\Z$ has type $(0,0,\dots)$ and $B$ has
type $(1,1,\dots)$. Both have endomorphism ring $\Z$, because neither type has an
entry $\infty$, while the types differ, so the groups are not isomorphic;
see~\cite{fuchs}. Rank $0$ is impossible when both groups are torsion, by the
Baer--Kaplansky theorem~\cite{fuchs}. We do not need these facts.
\end{remark}

\begin{thebibliography}{9}
\bibitem{fuchs} L.~Fuchs, \emph{Abelian Groups}, Springer Monographs in
  Mathematics, Springer, 2015.
\end{thebibliography}

\end{document}
